% =====================================================================
% Chapter 6 — Why Route A, Why Route B: The Two Routes Compared
% =====================================================================
\chapter{Why \routeA, Why \routeB: The Two Routes Compared}
\label{ch:comparison}

This chapter is the conceptual core of the dissertation. \secref{ch:routeA}
and \secref{ch:routeB} each justified their own route; this chapter
\emph{contrasts} them directly so that the choice is explicit and the
difference in output meaning is unmistakable. It first positions the two
routes as different answers to the same question under different data
availability (\secref{sec:cmp_positioning}), then compares them side by side
on motivation, assumptions, input requirements, ground truth and evaluation
(\secref{sec:cmp_table}), and finally discusses the \emph{output differences}
in depth (\secref{sec:cmp_outputs}).

\section{Positioning: two routes, one question}
\label{sec:cmp_positioning}

The research question of \secref{sec:research_question} --- how to turn
incomplete sonar returns into a useful 3D representation --- admits two
fundamentally different answers, and which one is \emph{correct} depends on
data availability and on what the output will be used for. The two routes are
\emph{complementary, not interchangeable}:

\begin{quote}
\routeA{} is the right answer \emph{when several scans of a rigid object
exist} and the output must be a \emph{fusion of measurements}: it estimates
poses, aligns the scans, and surfaces the measured evidence. \routeB{} is the
right answer \emph{when only one view exists} and a \emph{complete} object
representation is needed: it encodes the partial and decodes the missing
geometry from a learned prior.
\end{quote}

This is why the project implements both: neither route dominates the other,
because they answer different sub-problems of the same question. The decision
procedure is:

\begin{enumerate}
  \item \textbf{How many views of the object are available?} If several
  overlapping views exist and can be aligned, \routeA{} is available and its
  output is strictly more trustworthy (measured). If only one view exists,
  only \routeB{} can produce a complete output.
  \item \textbf{What is the output for?} If the application can tolerate (or
  requires) unmeasured regions, \routeB{}'s completeness is a feature. If the
  application must not be misled by hallucinated geometry (inspection, safety
  assessment), \routeA{}'s honesty is decisive.
  \item \textbf{Is a shape prior available and trusted?} \routeB{} depends on
  the training distribution: it will produce plausible-but-wrong geometry for
  shapes outside that distribution. \routeA{} needs no prior beyond object
  scale.
\end{enumerate}

\section{Side-by-side comparison}
\label{sec:cmp_table}

\tabref{tab:route_comparison} compares the two routes across every dimension
that a reader or examiner might question: motivation, assumption, input,
output, ground truth, failure modes and evaluation.

\begin{table}[H]
  \centering
  \small
  \begin{adjustbox}{max width=\textwidth}
  \begin{tabular}{@{}p{2.9cm}p{5.6cm}p{5.6cm}@{}}
    \toprule
    Dimension & \routeA{} (geometric fusion) & \routeB{} (learned completion) \\
    \midrule
    \textbf{Chosen when} & several scans of the same rigid object exist &
    only one partial view exists (or will exist at deployment) \\
    \textbf{Core assumption} & scans are noisy partials of one rigid object;
    there exist rigid transforms aligning them & training pairs teach a prior
    over complete object geometry; a complete reference exists for training \\
    \textbf{Input} & $N$ segmented 3~m scans (24 curated per object) &
    one segmented partial cloud (2048 points) \\
    \textbf{Internal operations} & FPFH $\to$ RANSAC $\to$ ICP $\to$ pose graph
    (MST + loop closures) $\to$ fusion $\to$ Poisson/BPA (or ray TSDF) &
    PointNet-style encoder $\to$ global latent $\to$ coarse head $\to$ folding
    decoder \\
    \textbf{Output} & fused measured cloud + two meshes; vertices with no
    measurement support are removed & coarse + fine completed point clouds;
    unobserved regions are inferred \\
    \textbf{Output meaning} & \emph{what was measured}, aligned and fused &
    \emph{what was inferred}, guided by a learned prior \\
    \textbf{Ground truth} & none external: the scans themselves are the
    evidence (no CAD mesh, no logged pose in \suop{}) & complete reference
    cloud paired with each partial (procedural or Route~A-mesh-derived) \\
    \textbf{Typical failure} & registration error propagates through the
    graph; symmetric parts confuse FPFH; sparse regions over- or under-filled
    by the mesher & hallucination: plausible-but-wrong geometry outside the
    training distribution; scale errors \\
    \textbf{Evaluation metric} & registration fitness/RMSE, graph
    connectivity, fused-point counts, mesh statistics, measured-to-surface
    residual & Chamfer distance (accuracy + completeness), visual
    partial/reference/completion comparison \\
    \textbf{Do not compare with} & \chd{} loss of \routeB{} & mesh counts of
    \routeA{} \\
    \bottomrule
  \end{tabular}
  \end{adjustbox}
  \caption{Direct comparison of the two routes. The critical column is
  ``Output meaning'': \routeA{} fuses measurements, \routeB{} infers missing
  geometry. Consequently their evaluation metrics measure different questions
  and must not be exchanged.}
  \label{tab:route_comparison}
\end{table}

\section{Output differences in depth}
\label{sec:cmp_outputs}

\subsection{What each output claims}
The most important intellectual point of this dissertation is that the two
outputs make \emph{different epistemological claims}:

\begin{itemize}
  \item \routeA{} claims: \emph{``every output point is a registered sonar
  return (or a surface interpolated between such returns); where no
  measurements exist, the surface is either left open (BPA) or explicitly
  trimmed (Poisson, ray TSDF).''} The output is an \emph{evidence report}.
  \item \routeB{} claims: \emph{``this is a complete object as predicted by a
  prior learned from the training distribution; some of these points were
  never measured and are inferred.''} The output is a \emph{model hypothesis}.
\end{itemize}

This distinction has three practical consequences.

\paragraph{1. Completeness versus fidelity.}
\routeA{} is typically less complete: parts of the object that no scan saw
remain empty, and the two meshers disagree there (Poisson closes, BPA stays
open). \routeB{} is always complete by construction: the decoder emits the
same fixed number of points regardless of how much of the object was visible.
Completeness is only meaningful relative to the ground truth, which \routeA{}
does not have in \suop{} and \routeB{} does.

\paragraph{2. Ground-truth availability.}
Because \routeA{}'s input is the only evidence, its accuracy cannot be checked
against an external truth in \suop{} --- there is no CAD model and no logged
pose. The honest metrics are therefore internal consistency (fitness, RMSE,
connectivity, fused counts) and measured-to-surface residuals, plus a
claim-free visual comparison. \routeB{} has an explicit paired reference, so
its accuracy \emph{can} be measured --- but only within the training
distribution, and a low Chamfer loss does not imply that the inference is
physically correct (the network may have learned the prior rather than the
specific object).

\paragraph{3. What ``better'' means.}
Because the two outputs answer different questions, a single ``which route is
better?'' is ill-posed. \routeB{}'s Chamfer loss is not comparable with
\routeA{}'s mesh statistics: a large \routeA{} mesh is not better than a small
one, and a small \routeB{} loss on synthetic partials is not evidence that the
geometric route failed. The only meaningful comparison would require a common
reference --- a calibrated acquisition with logged poses and an external scan
of the same objects --- which is exactly the future-work experiment of
\secref{ch:conclusions}.

\section{When the routes connect}
\label{sec:cmp_connection}

Despite their differences, the routes are not isolated. \routeB{}'s
Generator~2 (\secref{sec:routeB_generators}) samples training references from
\routeA{} meshes, so the learned prior is conditioned on what the geometric
route actually recovered from sonar. This creates a virtuous loop: where
\routeA{} has strong evidence, the prior is grounded in measured geometry; and
where \routeB{} then infers a complete shape, it does so from a distribution
that matches the sensor's output rather than an idealised CAD shape. The two
routes are best understood as partners: \routeA{} establishes the measured
core, \routeB{} extrapolates it into a complete model.

\section{Summary}
\label{sec:cmp_summary}

The choice between \routeA{} and \routeB{} is governed by data availability and
output purpose, not by which pipeline is ``more advanced''. \routeA{} answers
``what did we measure?'' with a pose-corrected fusion of the evidence;
\routeB{} answers ``what is the whole object?'' with a prior-informed
completion. Their outputs differ in completeness, in the status of unobserved
regions, and in the availability of ground truth --- and therefore their
evaluation protocols differ too. The next two chapters report the experimental
results of each route in turn.
